%\vspace{-0.1in}
\section{Introduction}
\label{section: introduction}

%\subsection{Background}
Recent advancements in large language models (LLMs) have highlighted their strengths in semantic understanding and contextual reasoning, enabled by extensive pre-training on vast corpora. Despite these strengths, LLMs often struggle with domain-specific queries and complex contexts, frequently generating ``hallucinations'', in which outputs appear credible but lack factual accuracy. 
Retrieval-augmented generation (RAG) addresses these limitations by integrating external knowledge to enhance factual accuracy and contextual relevance.
Early RAG implementations, however, often relied on document retrieval methods that introduced noise and excessive context~\cite{lewis2021retrieval,ColBERT,yu2023generate, HyDE}. 
This limitation has led to a growing focus on graph-based RAG systems, known as {\it GraphRAG}~\cite{GraphRAG,surveyTJFD,GRAG,SubgraphRAG}. 

In GraphRAG, conventional document retrieval is replaced by graph-based retrieval, leveraging the structured and relational nature of graphs to extract query-specific reasoning paths that provide more precise and contextually relevant support for LLM reasoning.
Typically, the GraphRAG workflow consists of two primary phases:
\begin{itemize}[leftmargin=12pt]
\item {\bf Retrieval:} This phase retrieves knowledge from graphs by identifying reasoning paths relevant to the query. 
\item {\bf Augmented generation:} The retrieved reasoning paths are used to augment the LLM prompt, enhancing its ability in reasoning and generating accurate and contextually relevant outputs. 
\end{itemize}
Existing GraphRAG studies~\cite{GraphRAG,GCR, RoG,LessIsMore,StructGPT,ref:kelp,ToG,DoG}
%, including Microsoft's RAG Project~\cite{GraphRAG}, 
have demonstrated the potential of integrating graph data management techniques with GraphRAG. However, despite their promise, these studies are still in their early stages and face two significant challenges. First, they lack foundational support in addressing the scalability of graph algorithms, which is crucial for managing large-scale graphs and reducing query latency in real-world applications. Second, there is a knowledge gap regarding the impact of semantic information on graphs on the performance of GraphRAG. Specifically, it is unclear which strategies are most effective for leveraging semantic information from both the query and the graph, and which type of semantic model is best suited for different query scenarios. Resolving these uncertainties is essential to improving the quality, efficiency, and cost of GraphRAG across various query scenarios.

%{\color{red}For instance, foundational support from database research is lacking, particularly in aspects such as efficient graph data storage and management, scalable graph query processing, and robust graph indexing---all vital for handling the dynamic and complex graph structures that GraphRAG depends on. {\bf Rewrite after findings are done.}}

%{\bf Is it possible to list a few key issues (paragraphs) when defining a GraphRAG, say scalability, accuracy, etc? If so, we can package them up into a few paragraphs, so that the motivation part can be better ``motivated''. Maybe, it can be done when experimental studies and key findings are finished.} {\color{red} Or just forget about it and briefly talk about connections with db research in the end of last paragraph.}

%\vspace{-5pt}

\textbf{Our Motivation.}
Building on the strengths of database research, we aim to address key challenges in GraphRAG and lay the groundwork for future advancements. 
While efforts like Modular RAG~\cite{modularRag} have modularized general RAG, they fall short of addressing the GraphRAG's workflows and needs.
%While prior effort (Modular RAG ~\cite{modularRag}) has made progress in modularizing general RAG workflows, it does not address the specific workflow and requirements of GraphRAG. 
%To this end, we identify three critical issues that should be addressed in order to move the field forward.
We identify three critical gaps:


\begin{itemize}  [leftmargin=12pt]
\item 
\textbf{Need for a Unified Framework for Categorizing and Analyzing GraphRAG Solutions.}
GraphRAG solutions consist of diverse technologies, including algorithmic approaches (e.g., random walk, PageRank) and neural network-based methods (e.g., end-to-end and LLM models), each playing a distinct role in their functionality. The lack of a unified categorization for systematically analyzing these techniques—despite prior efforts such as Modular RAG  ~\cite{modularRag} and other surveys ~\cite{RAGsurvey1,RAGsurvey2,RAGsurvey3,RAGsurvey4} on LLMs with graphs—hinders effective research summarization and the identification of promising candidates for further study.


\item 
\textbf{Need for Modular Retrieval in GraphRAG.}
Current implementations of GraphRAG often treat the retrieval phase as a single, monolithic process, making it difficult to isolate, analyze, and optimize individual components. A modular approach is desirable to facilitate the development of more efficient and effective retrieval mechanisms.
Moreover, while Modular RAG~\cite{modularRag} modularizes general RAG workflows, it neither dissects the core retrieval process of GraphRAG nor clarifies its distinction from general RAG.
%boundary from general RAG. 


\item 
\textbf{Absence of a GraphRAG Testbed for New Instance Design and Evaluation.}
Optimizing GraphRAG performance requires navigating a complex trade-off between runtime efficiency and reasoning accuracy, influenced by various design factors. A GraphRAG testbed is needed to generate diverse implementation instances, allowing users to evaluate, compare, and apply different approaches while providing clear guidelines and trade-offs for constructing optimal GraphRAG instances tailored to specific scenarios.
However, Modular RAG  ~\cite{modularRag} and related surveys ~\cite{RAGsurvey1,RAGsurvey2,RAGsurvey3,RAGsurvey4} focus on conceptual overviews and lack implementation or evaluation support for GraphRAG.
\end{itemize}

\begin{comment}


\subsection{Our Motivation}

Building on the strengths of database research, we aim to address key challenges in GraphRAG and lay the groundwork for future advancements. To this end, we identify three critical issues that should be addressed in order to move the field forward.


\begin{itemize}  [leftmargin=10pt]
\item \textbf{The Need for a Unified Framework for Categorizing and Analyzing GraphRAG Solutions.}
GraphRAG solutions consist of diverse technologies, including algorithmic approaches (e.g., random walk, PageRank) and neural network-based methods (e.g., end-to-end and LLM models), each playing a distinct role in its functionality. The lack of a unified categorization for systematically analyzing these techniques hinder effective research summarization and identification of promising candidates for further study.


\item \textbf{The Need for Modular Retrieval in GraphRAG.}
Current implementations of GraphRAG often treat the retrieval phase as a single, monolithic process, making it difficult to isolate, analyze, and optimize individual components. A modular approach is desirable to facilitate the development of more efficient and effective retrieval mechanisms.


\item \textbf{The Absence of a GraphRAG Testbed for New Instance Design and Evaluation.}
Optimizing GraphRAG performance requires navigating a complex trade-off between runtime efficiency and reasoning accuracy, which is influenced by various design factors. A GraphRAG testbed is needed to generate diverse implementation instances, allowing users to evaluate, compare, and apply different approaches while providing clear guidelines and trade-offs for constructing optimal GraphRAG instances tailored to specific scenarios.
\end{itemize}
\end{comment}

%}

%\textbf{Lack of a Unified Categorization for Analyzing and Organizing Existing GraphRAG Solutions.}The diverse components of a GraphRAG system, including algorithmic methods (e.g., random walk, PageRank) and neural network-based methods (e.g., end-to-end and LLM models), each play distinct roles in its functionality. The lack of a unified categorization for systematically analyzing and organizing these components hinders effective research summarization and identification of promising candidates for further study.
%However, the absence of a unified system to systematically categorize and compare these varied components impedes the effective summarization of existing research and the identification key trends in the design space and evolution of GraphRAG.
%This lack of a unified framework also impedes the subsequent stages of modularity design and performance evaluation.
%While many solutions (i.e., algorithms and neural network (NN) models) exist for constructing GraphRAG instances, the absence of a framework for systematically categorizing and comparing them limits the ability to summarize prior work and make informed choices for specific scenarios.

%\textbf{Limited Modularity in the Retrieval Phase of GraphRAG.} 
%Current GraphRAG implementations treat the retrieval phase as a single, unified process, making it difficult to isolate, analyze, and optimize individual components. By introducing a modular framework, we can clearly delineate the roles of each component, enabling targeted optimizations and system evaluations. %Moreover, a modular framework facilitates the creation of novel GraphRAG instances, unlocking new design possibilities and insights.

%obscuring the contributions of individual components and limiting the ability to analyze and optimize specific modules. With the unified framework, we can propose a modular design for the retrieval phase, allowing for clearer identification of distinct components and their roles. This approach not only enables targeted optimizations but also supports a more systematic and fair evaluation of GraphRAG implementations.
%Current GraphRAG instances treat retrieval as a monolithic process, obscuring the contributions of potential modules and hindering the analysis and optimization of specific ones.

%\textbf{No Testbed to Design New GraphRAG Instances.}Optimizing GraphRAG performance relies heavily on various design factors that influence the runtime efficiency and reasoning accuracy. Addressing these challenges necessitates a framework capable of generating diverse implementation instances for users to evaluate, compare, and apply, providing clear guidelines and trade-offs to construct optimal instances tailored to specific scenarios.
%GraphRAG performance depends on design factors related to solution efficiency and characteristics. However, without clear guidelines to illustrate the trade-offs between performance and these factors, practitioners encounter challenges in constructing optimal GraphRAG instances for specific scenarios.

%\vspace{-8pt}
\textbf{Our Contributions.}
We presents the first empirical study on GraphRAG and introduces {\it {\ourmethod}}, a unified and modular framework that provides insights to guide future research through contributions in both framework design and empirical analysis. 

For Framework Design:

\begin{itemize} [leftmargin=12pt]
    \item We propose LEGO-GraphRAG, a modular framework dividing the retrieval phase into two flexible modules: {\it {\preRetrieval}} and {\it {\retrieval}}, and classifies the techniques into {\it structure-based} and {\it semantic-augmented} methods for  each module.
    \item LEGO-GraphRAG’s modular framework, with categorized techniques, supports implementing all existing GraphRAG instances while enabling the creation of new ones, promoting both standardization and innovation.
    \item We identify essential design factors, including reasoning quality, efficiency, and cost, providing a structured trade-off analysis to guide the development of GraphRAG instances.
\end{itemize}
For Empirical Research:

\begin{itemize}[leftmargin=12pt]
    \item We study a comprehensive set of GraphRAG instances, including $7$ existing implementations and $16$ new instances generated by LEGO-GraphRAG. These instances were extensively evaluated on large-scale real graphs (i.e., Freebase) and $5$ commonly used GraphRAG query sets, covering various query scenarios.
    \item We suggest several modular configurations and a number of improvements to existing implementations for improved reasoning quality and balancing efficiency and cost.
\end{itemize}

%{\bf (4) Development of {\InstanceNum} High-Quality GraphRAG Instances.} Using {\ourmethod}, we construct {\InstanceNum} high-quality GraphRAG instances by combining representative algorithms and NN models, ensuring broad coverage of solution types across modules.

%{\bf (5) Comprehensive Empirical Analysis.} Through extensive studies on {\InstanceNum} instances, we analyze how different solution types impact retrieval and reasoning performance in LLM contexts, yielding many insights for effective GraphRAG development:

%The major findings of our study are the following:
%\subsubsection{Major Findings of Our Study}

%\begin{itemize}[leftmargin=15pt]
%	\item Findings1: GraphRAG involves a delicate trade-off between quality, efficiency, and cost, .
%	\item Findings2:XXX
%\end{itemize}

\subsection{Related Works}
\label{section: rw}

\subsubsection{RAG and GraphRAG}
Early text-based RAG systems rely on basic retrieval techniques, such as text chunking and cosine similarity for ranking~\cite{lewis2021retrieval}, which are prone to retrieving noisy and irrelevant information, leading to lower-quality results~\cite{HyDE}. To address this, recent works have 
introduced both pre- and post-retrieval improvements. 
Pre-retrieval methods~\cite{StepBackRAG} enhance the input query, 
while post-retrieval techniques~\cite{G-RAG, R4} refine the ranking and filtering of retrieved results. 
Toolkits like LlamaIndex and LangChain~\cite{LlamaIndex, LangChain} offer modular control over these stages, improving overall retrieval precision and system interpretability~\cite{HyDE}.
However, refining the retrieval stage to reliably suppress irrelevant content remains a key challenge~\cite{GraphRAG}. 
%The Modular RAG framework~\cite{modularRag} formalizes the general RAG pipelines but stops short of treating GraphRAG as a distinct, modular workflow. 

GraphRAG has recently emerged as a promising direction by representing knowledge in graph form, enabling more structured retrieval, multi-hop reasoning with  better interpretability~\cite{GraphRAG, RAGvsGraphRAG, 2501.13958}.
%Recently, GraphRAG has become a key focus in the RAG community, as representing knowledge as structured graphs enables finer-grained control, supports multi-hop reasoning with interpretability~\cite{GraphRAG, RAGvsGraphRAG, 2501.13958}. 
Compared to unstructured text retrieval, graph-based retrieval offers reduced noise, better coverage of entity relations,
%entity-level connections across text, 
and lowered token overhead during inference~\cite{ToG2.0, GraphRAG, NewArticles, Podcast}.\footnote{A detailed comparison between GraphRAG and text-based RAG is provided in our technical report B.10.} 
Existing instances (or implementations) in GraphRAG have drawn heavily from knowledge-base question answering (KBQA) techniques, utilizing both information retrieval methods~\cite{ref:wikimovie, EmbedKGQA} and semantic parsing models~\cite{QGG, HGNet} to identify relevant subgraphs or reasoning paths. 
Microsoft GraphRAG~\cite{GraphRAG} was an early effort exploring the advantages of graph-based over text-based retrieval, focusing on graph construction from text and precomputation. 
GCR~\cite{GCR} combines LLMs and beam search to refine reasoning paths, improving retrieval alignment with queries. RoG~\cite{RoG} and GSR~\cite{LessIsMore} employ Personalized PageRank (PPR) to retrieve query-specific subgraphs, while KELP~\cite{ref:kelp} refines reasoning paths using fine-tuned models like BERT~\cite{bert}. These studies highlight the diverse graph-based and neural network-based techniques employed in GraphRAG, showcasing the extensive potential of graphs in enhancing the reasoning quality of LLMs.

\subsubsection{LLMs with (Knowledge) Graphs}
GraphRAG aligns with the broader research to integrate LLMs with structured knowledge, such as KGs. Surveys~\cite{LLM_KG_Roadmap, LLM_KG_Trend} outline three paradigms: KG-enhanced LLMs, LLM-augmented KGs, and their joint integration.
%GraphRAG can be naturally situated within the broader research area of LLMs with (knowledge) graphs. This area has been discussed in recent surveys~\cite{LLM_KG_Roadmap, LLM_KG_Trend, SurveyKPLM, LLM_interface}, which review efforts to integrate LLMs and structured knowledge sources. \cite{LLM_KG_Roadmap} outlines three representative paradigms: KG-enhanced LLMs, LLM-augmented KGs, and synergized LLMs + KGs. 
Moreover, \cite{SurveyKPLM} reviews knowledge-enhanced pre-trained language models (e.g., LLMs) for language understanding and generation, and \cite{LLM_interface} examines LLMs as interfaces for data pipelines, highlighting their integration with KGs in AI systems.
%Our work aligns with this line of research by focusing on KG-enhanced LLMs and their synergistic use. Specifically, we propose a modular GraphRAG framework that supports the systematic design and evaluation of reasoning-enhanced LLMs with knowledge graphs.
Our work fits within the KG-enhanced LLMs paradigm and aims to contribute a modular GraphRAG framework to support systematic design and evaluation of reasoning-augmented LLMs with graphs.



\begin{comment}

In recent years, the development of large language models (LLMs) has garnered significant attention and driven advancements across multiple fields. Through extensive pre-training on vast textual datasets and their massive network parameters, LLMs have demonstrated impressive capabilities in semantic understanding and contextual reasoning. However, LLMs exhibit certain limitations when addressing domain-specific queries or handling complex contexts that require external knowledge. For instance, LLMs are prone to generating ``hallucinations'', wherein the information produced appears credible but is factually incorrect and misaligned with real-world knowledge. To address these limitations, retrieval-augmented generation (RAG) has emerged as an effective solution. This approach enhances the reasoning performance of LLMs by incorporating external knowledge during the reasoning process, thereby improving both the factual accuracy and contextual relevance of generated answers. Early RAG frameworks primarily relied on document retrieval, where relevant documents are retrieved from knowledge repositories based on the query context and integrated into the input provided to LLMs, enriching their understanding of real-world information. However, document-based retrieval methods often introduce noise and excessive context~\citep{HyDE}. As a result, the latest trend in the RAG community focuses on developing graph-based RAG frameworks: i.e., {\it GraphRAG}.


The GraphRAG process replaces documents with graphs~\footnote{In this paper, the term "graph" typically refers to complex graphs containing semantic information, with knowledge graphs being a representative example. Throughout the remainder of this paper, we use ``graph'' specifically to denote knowledge graphs.} in the retrieval process of RAG, enabling the extraction of specific ``reasoning paths'' from the graph based on the query context to provide to  LLMs for reasoning. Since graphs offer a structured representation of knowledge, the retrieved reasoning paths ensure that the information provided to LLMs is both relevant and precise. Additionally, the extensive range of search and ranking algorithms available for graphs helps filter out noisy information, addressing the redundancy issues often encountered in document-based retrieval methods. Existing works and ongoing development frameworks (e.g., Microsoft's ``GraphRAG'' system), have demonstrated the superiority and potential of GraphRAG compared to previous RAG methods.


%~\footnote{In this paper, we use "reasoning paths" as the format of prompt, following the previous works.\cite{RoG,GNN-RAG,ref:kelp}}



A {\it critical observation} is that GraphRAG represents an emerging domain in graph data management, particularly within the context of LLMs, and warrants increased attention from the database community. This assertion is grounded in two primary considerations: First, the database community can provide essential data management support for the efficient deployment of GraphRAG. Given that the foundational structure of GraphRAG is a ``Graph'', its storage, retrieval, and processing are areas extensively explored within the database field. A closer integration of these research efforts with GraphRAG is crucial for enhancing its capabilities. Second, the database community can contribute advanced solutions to the retrieval processes that are central to GraphRAG. Since GraphRAG’s technological core relies on ``graph-based retrieval algorithms and models'', a area of focus in database research, building upon these established studies enables GraphRAG to selectively adopt the most relevant techniques to improve its performance.




First, there is currently no unified framework to systematically categorize and evaluate the available solutions in GraphRAG (i.e., algorithms and neural network (NN) models). This results in the inability to effectively summarize and classify existing work on GraphRAG, while also impeding the clear identification of the practical efficacy of specific solutions within the GraphRAG process.
Second, current studies often treats GraphRAG as a monolithic process without performing a modular decomposition, which obscures the varying contributions of potential modules to overall performance. A more granular, modular GraphRAG framework would facilitate the analysis of trade-offs between module performance and solution choices, providing guidance for designing GraphRAG  instances that meet specific scenario requirements.


Thus, in this paper, we propose a unified and modular research framework for GraphRAG, named {\it {\ourmethod}}, which establishes four key standards to address the aforementioned challenges: {\bf 1) Modularization of GraphRAG}: {\ourmethod} divides the process of retrieving "reasoning paths" into three interconnected and flexible modules: {\it {\preRetrieval}}, {\it {\retrieval}}, and {\it {\postRetrieval}}; {\bf 2) Solutions for GraphRAG}:  {\ourmethod} systematically summarizes and classifies the algorithms or neural network (NN) models available for each module, facilitating a clear understanding of the potential design space for GraphRAG instances; {\bf 3) Design Factors of GraphRAG}: {\ourmethod} identifies two major factors that influence the design of GraphRAG instances, namely {\it Graph coupling} and {\it Computational Cost}, and analyzes how these factors impact the available solutions for each module; {\bf 4) Design trade-offs in GraphRAG}: {\ourmethod} provides a detailed analysis of the trade-offs between the execution performance of each module and the design factors, thereby offering valuable guidance for the future development of GraphRAG instances.




Furthermore, leveraging the {\ourmethod} framework, we build 20 high-quality GraphRAG instances by combining the most representative algorithms or NN models from various types of solutions while ensuring comprehensive coverage of different solution types for each module.
Through extensive empirical studies on these instances, we conducted a thorough analysis of the impact of various types of solutions on overall retrieval performance and reasoning performance based on LLMs.
Moreover, by synthesizing the experimental findings, we identified several critical insights into the development of GraphRAG instances from multiple analytical perspectives.
	内容...
\end{comment}
